% ---- pg-014 ----
用很大程度上取决于研究目标.在 1.2 节中我们列出少量实际问题,用以说明统计方法的选择与研究目标之间的联系.这些问题加上书中的例子将对多元方法在不同领域的适用性提供正确评价.

多元方法能最自然地发挥作用之处,便是下列科学研究目标:

\begin{enumerate}[label=\arabic*.,leftmargin=*,itemsep=2pt]
\item \mvbold{数据简化或结构简化}\quad 在不损失有价值信息的情况下尽可能简单地将被研究的现象描述出来.希望这样能使解释变得容易些.
\item \mvbold{分类与分组}\quad 根据所测量的特征将一些“类似的”对象或变量分组.另外,或许需要一些分类规则,以便将对象归入明确定义的各组.
\item \mvbold{变量间依赖性的研究}\quad 人们对变量间关系的本质感兴趣.是否所有变量都相互独立?还是有一个或多个变量依赖于其他变量? 如果是后者,那又是怎样依赖的?
\item \mvbold{预测}\quad 为了根据某些变量的观测值预测另一个或另一些变量的值,必须确定诸变量之间的关系.
\item \mvbold{假设的构造与检验}\quad 对以多元总体参数形式陈述的多种特殊统计假设进行检验.这样做可以验证某些假设或增强事先建立的信念.
\end{enumerate}

我们引用 F. H. C. 马里奥特(Marriott)的一段话为多元分析做一简要的综述(参见文献[19]).这段话是在讨论聚类分析时说的,不过我们觉得它适用于更多的方法.无论何时,在你尝试或阅读某种数据分析方法时,应该记住这一点.它可以使你保持正确的眼光,不致被理论中某些漂亮词句所左右.

\begin{quote}
要是所得结果与你所获的看法不一致,不要接受一个简单的逻辑解释,也不要在某种图形表示中清楚显露出来,因为这些结果可能是错的.数值方法并非魔术,招致失败的原因有很多.它们对数据解释来说不过是有用的帮手,不是能将大量数字自动转变成一组科学事实的碎肉灌肠机.
\end{quote}

\section{多元方法的应用}
\phantomsection\label{sec:1.2}

有关多元方法应用的出版物近年来有了惊人的增加.现在已很难像本书先前那些版本那样用简短的讨论来概括这些方法各式各样的实际应用了.然而,为了说明多元方法的实用价值,我们对来自若干学科的研究成果作一个简要介绍.所介绍的内容是依据前一节中给出的目标类别组织的.当然,我们的许多例子具有多面性,可适用于一个以上的类别.

\mvsub{数据缩减或简化}

\begin{itemize}[leftmargin=*,itemsep=2pt]
\item 用一些进行放射治疗的癌症患者的变量数据构造一个进行放射治疗的患者的简单测量法(参见练习 1.15).
\item 用许多国家和地区的径赛运动记录为男女运动员建立一个成绩标准(参见文献[8]和[22]).
\item 由高精度扫描仪获得的多光谱图像数据被简化成一种形式,可被看做是一个二维的海岸线的图像(参见文献[23]).
\item 利用与产量及蛋白质含量有关的几个变量的数据,建立一个选择母体的标准以改善下几代豆类植物(参见文献[13]).
\item 由职业仲裁人获得的数据所导出的策略相似性矩阵.通过这个矩阵,可帮助职业仲裁人确定评价解决争端所用策略的维数(参见文献[21]).
\end{itemize}
